\documentclass[11pt,twocolumn]{article}

% --- Core Packages ---
% --- Essential System Packages ---
\usepackage[utf8]{inputenc}       % Input encoding
\usepackage[T1]{fontenc}          % Font encoding for vector fonts
\usepackage{lmodern}              % Latin Modern clean vector font
\usepackage[margin=0.75in]{geometry} % Standard compact layout margins

% --- Math & Layout Utilities ---
\usepackage{amsmath,amssymb,amsfonts} % Advanced math typesetting
\usepackage{graphicx}             % For including graphics and charts
\usepackage{booktabs}             % Beautiful, professional table horizontal rules
\usepackage{lipsum}               % For placeholder text generation

% --- Hyperlinks & Links ---
\usepackage{url}
\usepackage[colorlinks=true, linkcolor=blue, citecolor=blue, urlcolor=blue, breaklinks=true]{hyperref}

% --- Algorithms ---
\usepackage{algorithm}
\usepackage{algpseudocode}        % TitleCase syntax matching \While and \EndWhile

% --- Citation Settings (APS Style) ---
% Sets superscript citations matching typical physical review journals
\usepackage[superscript,biblabel]{cite}


% --- Custom Commands ---
% --- Custom Typography and Shorthand Commands ---
\providecommand{\Require}{\item[\textbf{Require:}]}
\providecommand{\Ensure}{\item[\textbf{Ensure:}]}


\begin{document}

% --- Document Metadata ---
% --- Document Metadata ---
\title{\bf Clear and Robust Two-Column Scientific Layout}

\author{First A. Author$^{1,*}$ and Second B. Author$^{2}$}

\date{\small $^1$Department of Physics, University of Science, Country\\
$^2$Department of Machine Learning, Institute of Technology, Country\\
\today}


% --- Full-Width Title and Abstract ---
\twocolumn[
  \begin{@twocolumnfalse}
    \maketitle
    \begin{abstract}
        Your abstract goes here. It provides a concise summary of your core scientific methodology, empirical results, and broader impact. This template is built on the robust, standard \texttt{article} class configured for a clean, stable two-column physics layout that avoids internal system version crashes entirely.
        \vspace{25pt}
    \end{abstract}
  \end{@twocolumnfalse}
]

% --- Main Body ---

\section{Introduction}
Welcome to your rock-solid two-column layout. This setup uses standard, highly compatible LaTeX engines that compile natively on any MacTeX installation without crashing. References are automatically formatted to use clean superscript numbering styles like this~\cite{attention2017}, mimicking the American Physical Society (APS) layout perfectly.

\section{Methodology}
In quantitative scientific modeling, clear mathematical formulations are essential. You can write inline expressions like $E = mc^2$ or standalone field models using the standard equations:

\begin{equation}
    \mathcal{L}(\theta) = -\frac{1}{N} \sum_{i=1}^{N} \left[ y_i \log \sigma(f(x_i; \theta)) + (1 - y_i) \log (1 - \sigma(f(x_i; \theta))) \right].
    \label{eq:loss}
\end{equation}

Equation~\ref{eq:loss} fits perfectly within the bounds of a single column.

\section{Algorithms and Implementation}
Pseudocode can be presented cleanly inside your column rules using the modern, non-conflicting \texttt{algpseudocode} syntax:

\begin{algorithm}[H]
\caption{Stochastic Gradient Descent Update}\label{alg:sgd}
\begin{algorithmic}
\State \textbf{Require:} Learning rate $\eta$, initial parameters $\theta_0$, objective $\mathcal{L}$
\State \textbf{Ensure:} Optimized parameters $\theta$
\State $\theta \gets \theta_0$
\While{not converged}
    \State Sample mini-batch $\mathcal{B}$ from dataset
    \State Compute gradient $g \gets \nabla_\theta \mathcal{L}(\mathcal{B}; \theta)$
    \State $\theta \gets \theta - \eta \cdot g$
\EndWhile
\State \Return $\theta$
\end{algorithmic}
\end{algorithm}

\section{Results and Discussion}
\lipsum[1-2]

\subsection{Data and Parameters}
Tables use clean horizontal rules from the \texttt{booktabs} package (\texttt{\textbackslash toprule}, \texttt{\textbackslash midrule}, and \texttt{\textbackslash bottomrule}) instead of rigid vertical grid lines.

\begin{table}[h]
\centering
\caption{Key hyperparameters used during our empirical evaluation.}
\label{tab:parameters}
\begin{tabular}{lcr}
\toprule
Hyperparameter & Symbol & Value \\
\midrule
Learning Rate   & $\eta$ & $3 \times 10^{-4}$ \\
Batch Size      & $B$    & $256$ \\
Weight Decay    & $\lambda$ & $1 \times 10^{-5}$ \\
\bottomrule
\end{tabular}
\end{table}

\subsection{Visual Graphics}
Standard figures scale directly to your individual column width. Set widths explicitly to \texttt{\textbackslash columnwidth} to prevent layout bleeding:

\begin{figure}[h]
    \centering
    \includegraphics[width=\columnwidth]{example-image}
    \caption{This is a center-aligned figure scaling uniformly within the bounds of a single text column.}
    \label{fig:single_col}
\end{figure}

If you ever need a massive image or table that must stretch horizontally across \textit{both} columns, simply add an asterisk to the environment name, like this: \texttt{\textbackslash begin\{figure*\}} or \texttt{\textbackslash begin\{table*\}}.

\section{Conclusion}
This layout balances visual alignment with bulletproof compilation reliability. You can confidently expand this draft into a full publication manuscript~\cite{goodfellow2016}.

\section*{Acknowledgments}
We thank our colleagues for their helpful commentary on this manuscript. This work was supported by the Template Design Foundation under Grant No.~000-12345.

% --- References Block ---
\small
\bibliographystyle{unsrt}
\bibliography{bibliography/bibliography}

\end{document}
